\Needspace{12\baselineskip}
\originalchapter{噪声和射电谱的校准}
\originalsection{Johnson噪声和散粒噪声}
电压均值为零不表示电信号没有可测结构。阻值 $R$ 的电阻在温度 $T$ 下，一侧电压功率谱密度为 $S_V(f)=4\kb TR$，单位 $\mathrm{V^2/Hz}$；它要求 $hf\ll\kb T$ 的经典低频极限。测量链传递函数为 $H(f)$，输出方差
\[
\avg{V_{\rm out}^2}=4\kb TR\int_0^\infty|H(f)|^2\dd f+V_{\rm amp}^2.
\]
必须校准整个增益带宽积，不能拿最大增益乘一个任意带宽代替。放大器电压噪声、电流噪声、输入电容、导线电阻和温度不均匀都可能改变阻值依赖。

Junior Lab噪声模块通过不同电阻和温度比较输出均方值，区分热噪声与电路背景。若在摄氏温标上做线性外推，理想热噪声部分在绝对零度消失；放大器背景必须先作为独立项估计，否则外推截距改变。现代SI把 $h,e,\kb$ 的数值固定为精确值，这些沿用课程传统名称的实验现在更多体现单位实现和模型一致性检验，而非改变定义常数；依据BIPM定义常数表，见书末。

独立电荷 $e$ 以平均率 $I/e$ 到达。长度 $T$ 的电荷方差 $eIT$，对应一侧电流谱 $S_I=2eI$。若测量装置把电流转换成电压，也要积分相应跨阻增益。相关输运的Fano因子 $F$ 可使 $S_I=2eIF$，不能把任意导体都视为理想独立发射源。电源纹波和工频峰是窄带干扰，不是白噪声本底。

\originalsection{21厘米射电观测的测量链}
Junior Lab的21厘米模块用可控抛物面天线、低噪声放大器、混频和数字采样测太阳亮温和银河中性氢谱。中性氢基态超精细跃迁在约1420 MHz，沿视线的速度使谱线频率移动。观测系统记录功率谱，先要去除接收机本底、校准频率和天线温度。

在 $h\nu\ll\kb T$ 下，Rayleigh--Jeans强度 $I_\nu=2\kb T_b\nu^2/c^2$ 定义亮温。天线温度是天空亮温与天线波束的加权平均；源小于波束时有波束稀释，不能直接把天线温度当作太阳真实表面温度。增益和接收机温度可用已知噪声源或热冷负载求出，频带内增益曲线应与线上测量保持一致。

取非相对论径向速度 $v_r/c\simeq(\nu_0-\nu)/\nu_0$，需要说明参考系并改正地球自转、公转和太阳局部标准静止系运动。对银河平面内近圆轨道，
\[
v_r=\left[\frac{R_0V(R)}R-V_0\right]\sin l,
\]
其中 $l$ 为银经。内银河切点法用极端速度推算 $R=R_0|\sin l|$ 的转速；有限角分辨、非圆运动、光学深度和多距离气体会限制解释。单条谱线的速度不能唯一确定所有云的位置。

积分噪声的理想辐射计尺度为 $u_T\sim T_{\rm sys}/\sqrt{B\tau}$，独立带宽 $B$ 和积分时间 $\tau$ 降低白噪声。但较长积分还可能增加增益漂移、射频干扰和基线变化。实际天文图和源手册中的第三方星系图均未复制，本册没有生成假称天空实测的频谱。

\originalsection{功率谱的归一化与增益带宽}
本卷采用一侧谱：实信号 $V$ 的方差 $\avg{V^2}=\int_0^\infty S_V(f)\dd f$。两侧谱在正负频率各分一半，白噪声常数因子不同；引用 $4\kb TR$ 必须同时说明一侧定义。有限采样的谱还依赖采样率、窗口和FFT归一化，任意FFT模方不是自动具有 $\mathrm{V^2/Hz}$ 的谱密度。

输入短路测放大器电压噪声，多个电阻测电流噪声与热噪声的不同阻值依赖。最简单模型的输出谱为
\[S_{\rm out}=|H|^2[4\kb TR+S_{e_n}+R^2S_{i_n}],\]
电压/电流噪声相关时还有交叉项。高频输入电容使真实阻抗不再纯 $R$，必须用频率相关转移而非一个直流增益。

有效噪声带宽定义 $B_{\rm eff}=\int_0^\infty|H(f)|^2\dd f/|H|_{\max}^2$。一阶低通 $H=G/(1+if/f_c)$ 给 $B_{\rm eff}=\pi f_c/2$，不等于 $f_c$。这个解析例子说明为什么仅测3 dB点不足以假设矩形带宽。不同温度/电阻条件的拟合还共享此标定误差。

\originalsection{散粒噪声为何有因子二}
随机电流写作 $I(t)=e\sum_j\delta(t-t_j)$，事件率 $\lambda$。去均值后的协方差为 $\avg{\delta I(t)\delta I(t')}=e^2\lambda\delta(t-t')$，两侧频谱为 $e^2\lambda=eI$，一侧为 $2eI$。这里事件的独立性和电荷相同是关键；热平衡电阻噪声与有偏电流的散粒噪声不能随意互换。

若观察带宽内有工频和其谐波，应先查屏蔽、接地和供电；直接把这些峰一起积分再解释成白噪声常数，会系统高估目标量。若人为删频点，也要按实际 $H(f)$ 的积分修正带宽。报告中的噪声谱与对应积分必须用相同处理规则。

\originalsection{射电热冷负载、波束和频谱}
接收输出功率可写 $P=G\kb B(T_A+T_{
m rec})$。热/冷标准温度 $T_h,T_c$ 的比 $Y=P_h/P_c$ 给 $T_{
m rec}=(T_h-YT_c)/(Y-1)$，随后得到增益。若 $Y$ 接近1，误差会放大；标准负载与天空的阻抗/波束耦合不一致也产生系统偏差。

天空亮温与方向有关，$T_A=\int B_n(\Omega)T_b(\Omega)\dd\Omega/\int B_n\dd\Omega$。一个比主波束小很多的均匀圆源，其温度被面积填充因子稀释；旁瓣还会收入地面和银河背景。扫描太阳得到的角宽需扣除波束卷积，不能直接把观测宽度当太阳直径。

数字采样实信号的Nyquist频率为 $f_s/2$，频谱分辨约 $1/T_{
m record}$，窗函数改变主瓣宽和等效带宽。混频后的频率还可能反向：本振上/下边带应通过已知频率信号标定。基线扣除要避开H I线窗，过高阶多项式可以把宽线一起扣掉。

\originalsection{21厘米速度到银河转速的推导}
取太阳位置 $R_0$，天体圆轨道速度 $V(R)$，两者角速度差乘投影几何给 $v_r=[R_0V(R)/R-V_0]\sin l$。内银河视线切点为离银河中心最近处，$R_t=R_0|\sin l|$；若假设该处气体给极端径向速度，则 $V(R_t)=|v_{r,\max}|+V_0|\sin l|$（需按象限和符号一致写）。这不是对任意谱峰都适用的反演。

同一内银河 $R$ 常对应近、远两个距离；速度弥散、非圆棒结构和有限波束使切点也不完全尖锐。将频谱换为速度轴时还要说明采用射电或光学速度定义及参考系。本文只推导可测量链和条件，没有为银河绘制伪造观测转速曲线。
